% 原创教学几何示意：同向不同长度与等长不同方向。
\begin{tikzpicture}[figure picture,foundation picture,x=1mm,y=1mm,font=\figurefont\small]
  \node at (21,40) {端点距离};
  \node at (79,40) {方向夹角};
  \foreach \shift in {0,58}{
    \begin{scope}[xshift=\shift mm]
      \draw[foundation flow] (0,0)--(40,0) node[right] {$x_1$};
      \draw[foundation flow] (0,0)--(0,33) node[above] {$x_2$};
    \end{scope}
  }
  \draw[foundation flow,FigureRed,line width=1.1pt] (0,0)--(21,28) node[above right] {$\bm b$};
  \draw[foundation flow,FigureBlue,line width=1.1pt] (0,0)--(10.5,14) node[left] {$\bm a$};
  \draw[FigureStroke,line width=1.1pt,dashed] (13.5,13)--(24,27);
  \draw[FigureStroke] (12.2,14)--(14.8,12) (22.7,28)--(25.3,26);
  \node[text=FigureStroke,align=center] at (32,16) {长度\\不同};
  \begin{scope}[xshift=58mm]
    \draw[foundation flow,FigureBlue,line width=1.1pt] (0,0)--(28,7) node[right] {$\bm a$};
    \draw[foundation flow,FigureRed,line width=1.1pt] (0,0)--(14,25.24) node[above right] {$\bm b$};
    \draw[FigureStroke,line width=1.1pt] (9.70,2.43) arc[start angle=14.04,end angle=60.98,radius=10];
    \node[text=FigureStroke] at (14,10) {$\alpha$};
  \end{scope}
\end{tikzpicture}
